% Source printed pp. 511-512 / PDF pp. 519-520.
\subsection{Higgs 丛及其变体}\label{src:10.2.2-subsection}
Hitchin 在研究 Riemann 曲面上的自对偶 Yang--Mills 方程时引入了 Higgs 丛 \cite{Hit87}. 光滑连通射影曲线 $C$ 上的 \newterm{Higgs 丛}{Higgs bundle}是对 $(F,\phi)$, 其中 $F$ 是 $C$ 上的向量丛, 态射 $\phi:F\to F\otimes\Omega_C$ 称为 \newterm{Higgs 场}{Higgs field}. 若对每个 $\phi$-不变子丛 $F'\subsetneq F$ 有 $\mu(F')<\mu(F)$ (相应地, $\mu(F')\leq\mu(F)$), 则称 $(F,\phi)$ 稳定 (相应地, 半稳定). Nitsure 利用 GIT 构造了拟射影良模空间 $M_{r,d}^{H}(C)$, 参数化固定秩 $r$ 和次数 $d$ 的半稳定 Higgs 丛的 S-等价类 \cite{Nit91}. Faltings 在 \cite{Fal93} 中给出了 $M_{r,d}^{H}(C)$ 的不使用 GIT 的构造, 这启发了我们在 \S\bookref{9.6} 中的论证.

虽然模空间不紧, 但存在称为 \newterm{Hitchin 纤维化}{Hitchin fibration}的固有态射 $\pi:M_{r,d}^{H}(C)\to\mathbb A^D$, 其中 $D=\frac12\dim M_{r,d}^{H}(C)$, 它将 Higgs 丛 $(F,\phi)$ 映为 $\phi$ 的特征多项式的系数. Hitchin 纤维化的纤维 $\pi^{-1}(s)$ 是 Lagrange 子簇, 可与 $s\in\mathbb A^D$ 所对应的谱曲线 $C_s\subseteq T^*C$ 的 Picard 簇等同.

Narasimhan--Seshadri 对应 (注记~\bookref{9.1.19}) 将 $\pi_1(C)$ 的不可约酉表示与 $C$ 上的稳定向量丛等同; 它可推广为将 $\pi_1(C)$ 在 $\mathrm{GL}_r(\mathbb C)$ 中的不可约表示与 $C$ 上的稳定 Higgs 丛等同 \cite{Sim92}. 这一结果是``\newterm{非 Abel Hodge 对应}{non-abelian Hodge correspondence}''的一部分, 后者将 $\pi_1(C)$ 的表示 (等价地, 局部系统) 的 Betti 模空间 $M_{\mathrm B}$, 平坦联络的 de Rham 模空间 $M_{\mathrm{dR}}$, 以及次数为 $0$ 的 Higgs 丛的 Dolbeault 模空间 $M_{r,0}^{H}(C)$ 联系起来. 虽然这些空间的代数结构不同, 但有自然的微分同胚 $M_{\mathrm B}\cong M_{r,0}^{H}(C)$ 和同胚 $M_{\mathrm{dR}}\cong M_{r,0}^{H}(C)$. 微分同胚诱导上同调的同构 $\mathrm H^*(M_{\mathrm B},\mathbb Q)\cong\mathrm H^*(M_{r,0}^{H}(C),\mathbb Q)$, 而该上同调上的权滤过与 Hitchin 映射的反常 Leray 滤过之间存在令人惊讶的等价关系. 这就是 $\mathrm P=\mathrm W$ 猜想, 最近才得到解决 \cite{HMMS25,MS24}. 关于 Higgs 丛的模空间和 $\mathrm P=\mathrm W$ 猜想的更详细介绍, 推荐 \cite{Hei24,Hos24}.

除 Higgs 丛之外, 参数化带有附加结构的层的模空间还有许多变体, 例如:
\begin{itemize}
\item 带标点光滑曲线 $(C,p_i)$ 上的\newterm{抛物丛}{parabolic bundle}, 由 $C$ 上的向量丛 $F$ 以及每个标记点 $p_i$ 处的加权旗 $0=F_{i,0}\subseteq F_{i,1}\subseteq\cdots\subseteq F_{i,k}=F\otimes\kappa(p_i)$ 组成 \cite{MS80};
\item 对 $(F,A\to F)$, 由曲线上的向量丛 $F$ 和从固定层 $A$ 出发的映射 $A\to F$ 组成. 当 $X$ 是光滑曲线且 $A=\mathcal O_X$ 时, 有若干稳定性条件可定义固有模空间 \cite{BD91,Tha94}. 当 $X$ 是 Calabi--Yau 三维簇且 $A=\mathcal O_X$ 时, 若 $F$ 是 $1$ 维纯层且 $\operatorname{coker}(S)$ 为 $0$ 维, 则称对 $(F,s:\mathcal O_X\to F)$ 为 \newterm{Pandharipande--Thomas 稳定对}{Pandharipande--Thomas stable pair}. 它们构成固有模空间, 用于定义 PT 不变量;
\item 曲线上的全纯三元组 $(F_1,F_2,\phi)$, 由两个向量丛及它们之间的映射组成 \cite{BGPG04}; 以及
\item 相对于微分算子层 $\Lambda$ 的 $\Lambda$-模, 由带有 $\Lambda$-作用的凝聚层 $F$ 组成 \cite{Sim94a}, 其特殊情形包括 Higgs 丛和局部系统.
\end{itemize}
